
\subsubsection{Resultados numéricos de la calibración}
\label{sec:pie-rey-resultados}

Sustituyendo en la ley de propagación~\eqref{eq:sigma-final} las componentes
estimadas en las secciones anteriores se obtiene el balance de incertidumbres de la
indicación del pie de rey que se presenta en la
Tabla~\ref{tab:pie-rey-balance}.

\begin{table}[H]
\centering
\caption{Balance de incertidumbres de la indicación del pie de rey y contribuciones a la varianza combinada.}
\label{tab:pie-rey-balance}
\begin{tabular}{lrr}
\toprule
Contribución & $\sigma_i$ ($\mu$m) & Peso en la varianza (\%) \\
\midrule
Contacto parcial $E$ & 14.43 & 31.2 \\
Repetibilidad $R$ & 12.25 & 22.5 \\
Cambio de escala $S$ & 9.62 & 13.9 \\
Contacto lineal $L$ & 2.89 & 1.2 \\
Superficie completa $J$ & 0.00 & 0.0 \\
Cruce de interiores $K$ & 0.00 & 0.0 \\
Resolución $r$ & 14.43 & 31.2 \\
$\sigma(L_i)$ & 25.84 & 100.0 \\
$\sigma(L_{BP})$ & 0.1260 & 0.002 \\
$\sigma_c(E) = \sqrt{\sigma^2(L_i) + \sigma^2(L_{BP})}$ & 25.84 & — \\
$U(E) = 2\,\sigma_c(E)$ & 51.68 & — \\
\bottomrule
\end{tabular}
\end{table}


La varianza combinada de la indicación, conforme a la
ecuación~\eqref{eq:sigma-Li-total}, es la suma en cuadratura de las siete
contribuciones evaluadas:

\begin{equation}
    \sigma^2(L_i) = \sigma^2(L_{i_E}) + \sigma^2(L_{i_R}) + \sigma^2(L_{i_S})
    + \sigma^2(L_{i_L}) + \sigma^2(L_{i_J}) + \sigma^2(L_{i_K}) + \sigma^2(L_{i_r}),
    \label{eq:sigma-Li-simbolico}
\end{equation}

es decir

\begin{equation}
    \sigma^2(L_i) =
    \begin{aligned}
        &\left(\frac{E}{2\sqrt{3}}\right)^2
        + \left(\frac{s_R}{\sqrt{n}}\right)^2
        + \left(\frac{\max(S_{E-I}, S_{E-P})}{2\sqrt{3}}\right)^2
        + \left(\frac{L}{2\sqrt{3}}\right)^2 \\[4pt]
        &\quad + \left(\frac{J}{2\sqrt{3}}\right)^2
        + \left(\frac{K}{2\sqrt{3}}\right)^2
        + \left(\frac{d/m}{\sqrt{3}}\right)^2 .
    \end{aligned}
    \label{eq:sigma-Li-simbolico-exp}
\end{equation}

Sustituyendo los valores obtenidos para cada componente,

\begin{equation}
    \sigma^2(L_i) = (14.43)^2 + (12.25)^2 + (9.62)^2 + (2.89)^2 + 0 + 0 + (14.43)^2
    \ \mu\text{m}^2
    \label{eq:sigma-Li-num}
\end{equation}

\begin{equation}
    \therefore\ \sigma(L_i) \approx 25.84\ \mu\text{m}
    \label{eq:sigma-Li-num-valor}
\end{equation}

El error de indicación del pie de rey es prácticamente independiente del valor
patrón en las condiciones de trabajo: la diferencia de coeficientes de dilatación es
nula ($\delta\alpha = 0$, ambos de acero) y la diferencia de temperaturas entre el
bloque y el instrumento también lo es ($\delta T = 0$, ambos a $21.5\ ^\circ$C).
Evaluando la ecuación~\eqref{eq:sigma-cuad-expandida} con esos valores, los
coeficientes de sensibilidad $[1+\delta\alpha\,\Delta T_i]^2$ y
$[-1+\alpha_{BP}\,\delta T]^2$ toman el valor unidad, mientras que los términos
$[L_i\,\delta\alpha]^2$ y $[L_{BP}\,\delta T]^2$ se anulan por completo. Los dos
restantes, asociados a $\sigma^2(\delta\alpha)$ y a $\sigma^2(\alpha_{BP})$, no se
anulan pero resultan despreciables: el primero vale
$9.4\times10^{-8}\ \mu\text{m}^2$ en el extremo del alcance
(ecuación~\eqref{eq:termino-termico}) y el segundo es de orden $10^{-9}\ \mu\text{m}^2$,
frente a $\sigma^2(L_i) = 667.6\ \mu\text{m}^2$. Con esa aproximación, la
ecuación~\eqref{eq:sigma-cuad-expandida} se reduce a

\begin{equation}
    \sigma_c^2(E) = \sigma^2(L_i) + \sigma^2(L_{BP})
    \label{eq:sigma-c-num}
\end{equation}

Al desaparecer del modelo todos los términos que llevan la longitud como factor, la
varianza combinada deja de depender del valor nominal del punto de calibración. Si
se dejan explícitas las componentes térmicas de la
ecuación~\eqref{eq:sigma-cuad-expandida} como función de la longitud $L$, el modelo
completo conduce a la forma cerrada

\begin{equation}
    U(E,L) = \sqrt{8.343\times10^{-11}\,L^2 + 2670.4}\ \mu\text{m},
    \label{eq:U-vecL}
\end{equation}

con $L$ expresada en milímetros. El coeficiente del término cuadrático reúne el
aporte de los términos térmicos, que resulta despreciable: en el extremo del alcance
ensayado ($L = 250$ mm) contribuye con $9.4\times10^{-8}\ \mu\text{m}^2$ frente a los
$2670.4\ \mu\text{m}^2$ del término independiente, de modo que $U(E)$ permanece
constante en los cinco puntos de calibración.

Sustituyendo el valor de la ecuación~\eqref{eq:sigma-LBP-total-valor}:

\begin{equation}
    \sigma_c(E) = \sqrt{25.84^2 + 0.126^2}\ \mu\text{m}
    \approx 25.84\ \mu\text{m}
    \label{eq:sigma-c-num-valor}
\end{equation}

Finalmente, la incertidumbre expandida con factor de cobertura $k = 2$ resulta

\begin{equation}
    U(E) = 2 \times 25.84\ \mu\text{m} \approx 51.68\ \mu\text{m}
    \label{eq:U-num-valor}
\end{equation}

\subsubsection{Declaración de conformidad}

El error máximo permitido del pie de rey 530-115 especificado por el fabricante es
$\text{EMPE} = \pm 80\ \mu\text{m}$ para todo su alcance de $0$ a $300$ mm, aunque
el ensayo cubrió de $50$ a $250$ mm. Dado que

\begin{equation}
    U(E) \approx 51.68\ \mu\text{m} < \text{EMPE} = 80\ \mu\text{m},
    \qquad \frac{U(E)}{\text{EMPE}} \approx 0.65,
    \label{eq:conformidad}
\end{equation}

la incertidumbre de la calibración es menor que la tolerancia del instrumento, por
lo que el resultado permite decidir conformidad. Los cinco errores de indicación se
mantienen por debajo del EMPE, con un máximo de $+28.2\ \mu\text{m}$ en $150$ mm, el
$35$\,\% del límite permitido.

La conformidad es, por tanto, ajustada. Adoptando la banda de aceptación de
ILAC-G8, que exige $|E| \le \text{EMPE} - U \approx 28.3\ \mu\text{m}$, tres de los
cinco puntos caen holgadamente dentro, mientras que los de $150$ mm y $200$ mm
quedan en el borde de la zona de incertidumbre, donde el error está dentro de
tolerancia pero la incertidumbre del método impide afirmar conformidad sin riesgo.
El instrumento cumple, pero el margen disponible frente al EMPE es reducido y
conviene repetir la calibración si el instrumento sufre golpes o desgaste
apreciable.
